\section{Анализ предметной области}
\subsection{Характеристика сферы поиска вакансий и подбора IT-специалистов}

В практике IT-рекрутинга поиск работы устроен иначе. Помимо должности и опыта здесь критичны конкретный набор технологий, формат занятости и уровень специализации. Во многих других сферах достаточно указать профессию, стаж и общие обязанности — в IT этого часто мало. На практике две вакансии backend-разработчика могут требовать совершенно разных навыков: где-то нужен Node.js и PostgreSQL, а где-то — Java, Spring и опыт работы с микросервисами.

Поэтому описание вакансии в IT должно быть точным и структурированным. Работодателю необходимо перечислить стек, требования, тип занятости и особенности задач — иначе вакансия может остаться незамеченной. Соискателю, в свою очередь, нужно сразу понять, подходит ли вакансия под его опыт и навыки, не тратя время на просмотр нерелевантных вариантов.

IT-сфера отличается высокой степенью специализации. Backend-разработчик, frontend-разработчик, тестировщик, DevOps-инженер, аналитик, fullstack-разработчик и системный администратор — все они из IT, но задачи и инструменты у них разные. Даже внутри одной специализации требования могут существенно различаться — всё зависит от стека, архитектуры проекта и организации работы в команде.

Главная особенность подбора IT-специалистов в том, что работодатель ищет не просто должность, а кандидата, чей набор навыков отвечает конкретным требованиям проекта. Соискатель, в свою очередь, оценивает зарплату, содержание задач, возможность удалённой работы и перспективы роста. Таким образом, процесс поиска строится на сравнении требований вакансии и компетенций кандидата, которое каждый участник выполняет самостоятельно, опираясь на предоставленные системой данные.

Не менее важен формат занятости. Для IT характерны удалённая работа, гибрид и офис, а также разные типы трудоустройства: полная, частичная занятость, стажировка или проектная работа. Это напрямую влияет на выбор: начинающий специалист ищет стажировку или junior-позицию, а опытный разработчик обращает внимание на уровень ответственности и условия. Сама специфика IT-рекрутинга требует, чтобы эти параметры были отражены в структуре вакансии и учитывались соискателем при поиске.

Таким образом, ключевые особенности IT-сферы — высокая специализация, значимость конкретного стека технологий и навыков, разнообразие форматов занятости и необходимость точного описания вакансий — определяют, какие данные должны быть представлены в вакансии и профиле кандидата. В проектируемой системе это отражается через следующие поля: специализация (specialization), навыки (skills), формат работы (work\_format), тип занятости (employment\_type), зарплата (salary) и местоположение (location).

\subsection{Роли пользователей и процессы взаимодействия в системе IT-рекрутинга}

Основные участники платформы IT-рекрутинга --- соискатель, работодатель (или рекрутер) и администратор. Каждый выполняет собственные действия и работает с разным набором данных, поэтому важно заранее определить их роли, права и типовые сценарии.

После регистрации соискатель заполняет профиль, который фактически выполняет роль резюме. Он указывает специализацию, навыки, ссылки на портфолио (например, GitHub). Это особенно важно в IT, где работодатель оценивает не просто должность, а техническое соответствие вакансии. Затем соискатель ищет вакансии, сравнивает условия и отправляет отклик на подходящие варианты. Также он может сохранять интересные вакансии и возвращаться к ним позже.

Работодатель решает связанную задачу: ему необходимо найти подходящего специалиста под потребность компании. Для этого он должен представить компанию, сформулировать вакансию, указать требования к кандидату и опубликовать предложение. В отличие от простого размещения объявления, работа работодателя не заканчивается после публикации вакансии. После получения откликов начинается следующий этап: просмотр профилей кандидатов, оценка их соответствия требованиям и изменение статуса отклика.

Процесс строится последовательно: работодатель публикует вакансию, соискатель изучает её и отправляет отклик, после чего работодатель рассматривает кандидата и меняет статус отклика. Статус (например, «новый», «на рассмотрении», «принят», «отклонён») показывает, на каком этапе находится заявка. Для соискателя это снижает неопределённость, а работодателю помогает поддерживать порядок при работе с несколькими кандидатами.

Администратор не участвует в подборе персонала и не принимает решений по кандидатам. Его роль ограничена базовым контролем: он просматривает списки пользователей и вакансий, блокирует пользователей, удаляет или снимает с публикации некорректные вакансии. Такой набор функций позволяет поддерживать порядок в системе без вмешательства в сам процесс рекрутинга.

Чтобы платформа работала корректно, права пользователей разделены. Соискатель ищет вакансии, управляет профилем и отправляет отклики. Работодатель создаёт вакансии и обрабатывает отклики кандидатов. Администратор выполняет только указанные выше контрольные действия. Такое разграничение снижает риск некорректного доступа к чужим данным и делает систему более управляемой.

Таким образом, распределение ролей выстраивает процесс IT-рекрутинга последовательно: от публикации вакансии до обработки откликов и контроля размещаемой информации.

\subsection{Проблематика существующих подходов к поиску работы и подбору персонала в IT-сфере}

На практике поиск работы и подбор персонала в IT-сфере часто ведутся с использованием множества не связанных друг с другом каналов: крупных сайтов вакансий, профессиональных сообществ, мессенджеров, социальных сетей, личных рекомендаций и внутренних таблиц работодателя. Каждый из этих инструментов может быть полезен по отдельности, однако их одновременное применение приводит к разрозненности данных о вакансиях, кандидатах и откликах. В таких условиях сложнее контролировать и модерировать информацию, из-за чего возрастает риск появления устаревших или недостоверных сведений. Высокая степень специализации, присущая IT-сфере, лишь усиливает негативные последствия этой разрозненности: неточное описание вакансии или несоответствие навыков оборачиваются ещё бо́льшими временны́ми потерями для обеих сторон.

Наиболее заметная трудность --- большое количество нерелевантных вакансий и откликов. Соискатель может находить предложения, которые формально совпадают по названию должности, но не соответствуют его реальному стеку технологий, уровню подготовки или желаемому формату работы. Например, разработчик, специализирующийся на backend-разработке с Node.js, получает в выдаче вакансии, связанные с Java, PHP или смежными направлениями. Работодатель сталкивается с обратной ситуацией: на опубликованную вакансию откликаются кандидаты, чьи навыки не отвечают ключевым требованиям, из-за чего время, затрачиваемое на первичный отбор, значительно увеличивается.

Эта проблема особенно заметна в описаниях вакансий. В IT одинаковые названия должностей далеко не всегда подразумевают одинаковый набор задач. Вакансия <<frontend-разработчик>> может предполагать разработку интерфейсов на React, поддержку существующего проекта на Vue или выполнение задач, близких к fullstack-направлению. Если требования не детализированы, соискатель вынужден тратить время на изучение вариантов, которые заведомо не соответствуют его профилю, а работодатель получает отклики от специалистов с неподходящим опытом.

Отсутствие прозрачного состояния отклика дополнительно усложняет взаимодействие. Соискатель, отправив значительное количество откликов, часто не имеет чёткого представления о том, был ли его отклик просмотрен, находится ли он на рассмотрении или уже отклонён. Когда обратная связь осуществляется через разрозненные каналы --- почту, мессенджеры, --- история взаимодействия легко теряется: сообщения затериваются, и ни соискатель, ни работодатель не видят целостной картины по каждому кандидату. Для работодателя ситуация оказывается не менее сложной: при большом потоке кандидатов трудно фиксировать, кто уже рассмотрен, кому необходимо ответить, а какие заявки требуют дальнейших действий.

Отдельную трудность для работодателя представляет управление вакансиями. Когда информация о вакансиях хранится в разных источниках, контролировать их актуальность становится крайне затруднительно. Закрытая вакансия может продолжать отображаться для соискателей; изменившиеся требования не обновлены своевременно; часть откликов остаётся без обработки. Это не только снижает удобство работы самого работодателя, но и подрывает доверие соискателей к достоверности публикуемых сведений.

Наконец, важной проблемой является отсутствие системной модерации. Поскольку вакансии и профили создаются самими пользователями, в неконтролируемой среде неизбежно появляются ошибочные, устаревшие или заведомо недостоверные данные. Без административного контроля невозможно своевременно выявлять и устранять некорректные вакансии, блокировать недобросовестных участников и поддерживать качество информации. В результате доверие к сервису снижается, а эффективность подбора падает.

Таким образом, ключевые трудности порождены разрозненным использованием множества каналов и отсутствием единой информационной среды, в которой данные о вакансиях, профилях кандидатов и статусах откликов были бы согласованы и защищены от искажений. Это замедляет процесс подбора и создаёт риски для всех участников. Указанные проблемы определяют необходимость перехода к специализированной платформе, реализующей базовый цикл взаимодействия между соискателем и работодателем в рамках единой системы.

\subsection{Обоснование необходимости разработки web-платформы}

Как показано в разделе 1.3, ключевые трудности, возникающие при разрозненном использовании каналов IT-рекрутинга, вызваны отсутствием единого информационного пространства и, как следствие, низкой прозрачностью и управляемостью процессов. Устранение этих проблем требует перехода от набора несвязанных инструментов к централизованной web-платформе, в которой базовый цикл взаимодействия между соискателем и работодателем реализован в рамках одной системы.

Единая платформа позволяет консолидировать публикацию вакансий, хранение профилей соискателей и данных компаний, отправку и обработку откликов в централизованной цифровой среде. Исключается дублирование данных, снижается вероятность потери информации, а действия всех участников становятся согласованными.

Принципиальное значение приобретает прозрачность состояний откликов. Платформа даёт возможность присваивать каждому отклику чётко определённый статус --- «новый», «на рассмотрении», «принят», «отклонён», --- что обеспечивает как соискателю, так и работодателю ясную картину движения заявки. Работодатель получает системный инструмент для параллельной работы с несколькими кандидатами, а соискатель получает информацию о том, как продвигается рассмотрение его обращения.

Чёткое разграничение прав доступа является ещё одним необходимым условием надёжного функционирования платформы. Соискатель, работодатель и администратор выполняют разные функции, поэтому их действия должны быть ограничены соответствующими разрешениями. Это предотвращает пересечение ролей и риск случайного или несанкционированного изменения чужих данных.

Для поддержания актуальности информации платформа должна обеспечивать управление жизненным циклом вакансий через набор статусов: «черновик», «опубликована», «снята с публикации». Благодаря этому работодатель может своевременно обновлять описание, снимать устаревшие предложения и тем самым избавлять соискателей от просмотра неактуальных вариантов.

В платформе также предусматривается базовый административный контроль, ограниченный функциями, необходимыми для поддержания порядка: просмотр списков пользователей и вакансий, блокировка нарушителей, удаление или снятие с публикации некорректных вакансий. Такой набор полномочий позволяет сохранять работоспособность системы, не наделяя администратора избыточными правами.

Учёт специфики IT-рекрутинга, подробно охарактеризованной в разделе 1.1, требует, чтобы поиск и фильтрация были не просто дополнением, а центральным инструментом работы с данными. Высокая степень специализации, значимость конкретного стека технологий и разнообразие форматов занятости делают необходимыми точные фильтры по навыкам, специализации, формату работы и типу занятости. Платформа предоставляет эти инструменты для самостоятельного поиска, позволяя соискателю отбирать вакансии, соответствующие его опыту и предпочтениям, а работодателю --- просматривать отклики и профили кандидатов без автоматических рекомендаций.

Разрабатываемая система является учебным проектом и имеет чётко очерченные границы минимально жизнеспособного продукта (MVP). За его пределами остаются: чат, email-уведомления, PDF-резюме, рейтинги компаний и кандидатов, избранные кандидаты, восстановление пароля, мобильное приложение и рекомендательная система. Сознательное ограничение функциональности позволяет сосредоточить усилия на базовом, наиболее значимом цикле взаимодействия между соискателем и работодателем.

Таким образом, разработка web-платформы для поиска работы и подбора IT-специалистов обоснована. Она даёт возможность объединить основные действия в рамках MVP в единой системе, обеспечить прозрачность и управляемость каждого этапа, предоставить участникам точные инструменты поиска и фильтрации, не выходя за пределы реализуемого в рамках дипломного проекта объёма функциональности.

\subsection{Структура предметной области и ключевые сущности системы}

На основании рассмотренной предметной области можно выделить основные сущности, определяющие структуру онлайн-платформы для поиска работы и подбора IT-специалистов. Эти сущности отражают как участников процесса IT-рекрутинга, так и данные, с которыми они работают в рамках системы.

Базовая сущность --- пользователь. Через учётную запись он получает доступ к функциям платформы, а его возможности определяются ролью. Роль не выделяется в отдельную сущность, а является ключевым атрибутом пользователя. В системе предусмотрены три основные роли: соискатель, работодатель и администратор.

Профиль соискателя содержит его профессиональные данные и выполняет функцию резюме. В нём указываются специализация, краткое описание, навыки и дополнительные сведения (например, ссылка на GitHub). Работодатель использует профиль для первичной оценки кандидата и сопоставления его компетенций с требованиями вакансии. Отказ от отдельной сущности <<Резюме>> исключает дублирование данных.

Для представления работодателя в системе используется самостоятельная сущность <<Компания>>. Она содержит сведения об организации, от имени которой публикуются вакансии и ведётся подбор специалистов, и напрямую связана с пользователем, имеющим роль работодателя. Компания не является профилем работодателя, а представляет отдельную сторону в процессе рекрутинга. Вакансии публикуются от имени компании, что делает её видимой для соискателей и позволяет связывать все предложения одного работодателя воедино.

Центральное место занимает вакансия. Это предложение работы, связанное с компанией и содержащее сведения о специализации, требованиях, условиях, формате занятости и состоянии публикации. Такие характеристики, как тип занятости и формат работы, не вынесены в отдельные справочники, а задаются непосредственно в вакансии, что упрощает структуру системы. Вакансия является основой взаимодействия между работодателем и соискателем, поскольку через неё начинается процесс отклика и последующего рассмотрения кандидата.

Важную роль играет навык. Он используется и в профиле соискателя, и в требованиях вакансии. Связь <<многие ко многим>> между профилем соискателя и навыками, а также между вакансией и навыками реализуется через отдельные ассоциативные сущности <<Навыки соискателя>> и <<Навыки вакансии>>. Это даёт кандидату возможность свободно указывать свои навыки, а работодателю --- перечислять несколько требуемых технологий в одной вакансии.

Связь между кандидатом и вакансией фиксируется сущностью <<Отклик>>. Она фиксирует факт обращения соискателя к конкретной вакансии и позволяет отслеживать дальнейшее продвижение заявки. Отклик является не просто отметкой о проявленном интересе, а отдельным элементом процесса подбора, поскольку через него работодатель рассматривает кандидата и принимает решение. При этом на уровне предметной области накладывается ограничение: один соискатель не может откликнуться на одну и ту же вакансию повторно.

Связь между пользователем и интересующими его вакансиями реализуется через сущность <<Избранные вакансии>>. Она позволяет сохранять вакансии для последующего просмотра без дублирования данных и не подразумевает функциональности избранных кандидатов.

Важное значение имеют статусы вакансии и отклика. Статус вакансии определяет текущее состояние предложения работодателя (черновик, опубликована, снята с публикации), а статус отклика отражает этап его рассмотрения (новый, на рассмотрении, принят, отклонён). Без этих атрибутов сложно обеспечить последовательную и прозрачную работу платформы.

Таким образом, структура предметной области включает основные сущности: пользователь, профиль соискателя, компания, вакансия, навык, отклик и избранные вакансии. Связи <<многие ко многим>> реализуются через ассоциативные сущности <<Навыки соискателя>> и <<Навыки вакансии>>, а важную роль играют атрибуты и справочные значения, прежде всего роль пользователя и статусы вакансий и откликов. Совокупность этих элементов отражает логику работы платформы и создаёт основу для последующего проектирования базы данных и программной логики системы.

\subsection{Анализ существующих решений в области IT-рекрутинга}

Для определения требований к разрабатываемой web-платформе рассмотрены существующие решения, применяемые для поиска работы и IT-рекрутинга. Анализ ограничен функциями, связанными с логикой проектируемой системы: разделение пользователей по ролям, профиль соискателя, профиль компании, публикация вакансий, указание навыков, поиск и фильтрация, отправка откликов, обработка статусов и сохранение избранных вакансий.

В качестве аналогов выбраны HeadHunter, Хабр Карьера и LinkedIn --- платформы, различающиеся масштабом, аудиторией и специализацией.

\subsubsection{HeadHunter}

HeadHunter предоставляет детализированную систему статусов откликов, включающую, например, «отправлен», «просмотрен», «отказ», «приглашение» (приведены обобщённые названия). Для проектируемой платформы данная модель упрощена до четырёх статусов: «новый», «на рассмотрении», «принят», «отклонён». Такого набора достаточно, чтобы отразить основные этапы рассмотрения кандидата и избежать избыточности, свойственной коммерческому сервису. HeadHunter также реализует базовый цикл «вакансия --- отклик --- рассмотрение», который положен в основу логики разрабатываемой системы.

Рекомендательная система, платные продвижения вакансий, расширенная аналитика и инструменты массового рекрутинга, присутствующие в HeadHunter, являются избыточными для целей учебного проекта и в рамках данной работы не рассматриваются.

\subsubsection{Хабр Карьера}

Хабр Карьера ориентирована на IT-сферу и группирует вакансии по специализациям (разработка, тестирование, администрирование и др.). Навыки выступают ключевой характеристикой как профиля соискателя, так и вакансии; поиск кандидатов ведётся через раздел «Специалисты» с фильтрацией по профессиональным компетенциям. В проекте данный подход отражён возможностью указания навыков и специализаций, а также возможностью поиска соискателей по этим параметрам. Кроме того, каждая вакансия привязывается к профилю компании --- принцип, реализованный в Хабр Карьере.

Рейтинги компаний, оценки работодателей, система подписок и платные инструменты продвижения, представленные на платформе, исключены из рассмотрения как функции, избыточные для учебной платформы.

\subsubsection{LinkedIn}

LinkedIn демонстрирует тесную интеграцию профиля специалиста, страницы компании и публикуемых вакансий. Это позволяет работодателю не только получать отклики, но и формировать доверительный образ, а соискателю --- видеть контекст работодателя. Инструмент активного поиска LinkedIn Recruiter поддерживает фильтрацию кандидатов по должностям, навыкам и отраслям. В учебной платформе данная идея реализована в упрощённом виде: работодатель может фильтровать профили соискателей, однако без сложных аналитических и рекомендательных механизмов.

Социальная составляющая LinkedIn (лента публикаций, сеть контактов, личные сообщения), а также расширенные карьерные страницы и мобильное приложение в рамках учебного проекта не предусмотрены, так как выходят за пределы минимально необходимого функционала.

\subsubsection{Вывод по анализу аналогов}

Все рассмотренные платформы подтверждают востребованность следующих элементов: профиль соискателя, профиль компании, публикация вакансий, указание навыков, поиск и фильтрация кандидатов и вакансий, отклики и управление их статусами.

Каждый аналог внёс характерный вклад в проектные решения:
\begin{itemize}
	\item HeadHunter --- развитая статусная модель откликов и понятный цикл обработки кандидатов;
	\item Хабр Карьера --- ориентация на IT-специализации и использование навыков как ключевого критерия поиска;
	\item LinkedIn --- связка «специалист --- компания --- вакансия» и инструменты активного поиска кандидатов.
\end{itemize}

На основе проведённого анализа выделены функции, целесообразные для включения в MVP:
\begin{itemize}
	\item регистрация и авторизация пользователей;
	\item профиль соискателя и компания работодателя;
	\item публикация вакансий с указанием навыков;
	\item поиск и фильтрация вакансий и кандидатов;
	\item отклики на вакансии и четыре статуса отклика: «новый», «на рассмотрении», «принят», «отклонён»;
	\item сохранение избранных вакансий;
	\item базовые административные функции (просмотр списков пользователей и вакансий, блокировка пользователей, снятие некорректных вакансий с публикации).
\end{itemize}

Функции, выходящие за рамки минимального учебного объёма (чат, email-уведомления, рекомендации, рейтинги, загрузка PDF-резюме, избранные кандидаты, мобильное приложение), на данном этапе не рассматриваются. Такой набор покрывает ключевые сценарии IT-рекрутинга: соискатель создаёт профиль и откликается на вакансии, работодатель публикует вакансии и управляет откликами, администратор контролирует пользователей и размещаемую информацию. Исключённые возможности могут быть добавлены при дальнейшем развитии системы, а их реализация в рамках дипломного проектирования нецелесообразна.



